\section{Experiments and Results}

% TODO

%* weak vs medium PersonalAI
%* medium PersonalAI vs GraphRAG analogues
%* max clue-queries ablation
%* plan enhancment disable
%* Latency
%* error analysis

%- сказать то достигнут sota результат по сравнению с таким-то набором методов
%- сказать, что перегенрация плана поиска даёт значительный прирост эффектинвновсти + дать ссылку на пример из аппендикса
%- обратить внимание на то, что использование алгоритмов обхода графа для поиска релевантной инофрмации даёт больший прирост эффективности по сравнению с flatten ретривом. Обратить внимание что добавление механизма планирования немного уменьшает разницу в эффективности
%- сказать, что увеличение числа clue-запросов повышает точность ответа из-за того, что при добавлении информации в память выполняется только посимвольная проверка + дать ссылку на пример из аппендикса
%- сказать про типы ошибок, которые возникают в PAI-2 + дать ссылки на примеры из аппендикса. сформулировать план действий для их устранения
%- сказать про скорось работы PAI. сказать, что время работы увеличилось, так как добавлен дополнительный механизм.

Based on conducted experiments, a comparative table, summarizing best-performing QA configurations by LLM-as-a-Judge metric, was compiled (see Table~\ref{tab:main_comparison}).

\begin{table*}[th!]
    \centering
	\renewcommand{\arraystretch}{1.3}
	\resizebox{\textwidth}{!}{%
    \begin{tabular}{|cc|c|c|c|c|c|c|c||c|}
    \hhline{|--|-|------||-|}
    \multicolumn{2}{|c|}{\multirow{2}{*}{\textbf{Method}}} & \multirow{2}{*}{\textbf{LLM}} & \multicolumn{6}{c||}{\textbf{Dataset}} & \multirow{2}{*}{Mean} \\ 
    \hhline{|~~|~|------||~|}
    \multicolumn{2}{|c|}{} &  & \multicolumn{1}{c|}{\textbf{Natural Questions}} & \multicolumn{1}{c|}{\textbf{TriviaQA}} & \multicolumn{1}{c|}{\textbf{HotpotQA}} & \multicolumn{1}{c|}{\textbf{2WikiMultihopQA}} & \multicolumn{1}{c|}{\textbf{MuSiQue}} & \textbf{DiaASQ} &  \\ 
    \hhline{=========::=}
    \multicolumn{2}{|c|}{LightRAG} & \multirow{15}{*}{Qwen2.5 7B} & \multicolumn{1}{c|}{0.26} & \multicolumn{1}{c|}{0.37} & \multicolumn{1}{c|}{0.15} & \multicolumn{1}{c|}{0.11} & \multicolumn{1}{c|}{0.01} & 0.07 & 0.16 \\ 
    \hhline{|--|~|-|-|-|-|-|-||-|}
    \multicolumn{2}{|c|}{RAPTOR} &  & \multicolumn{1}{c|}{0.66} & \multicolumn{1}{c|}{0.73} & \multicolumn{1}{c|}{0.46} & \multicolumn{1}{c|}{0.27} & \multicolumn{1}{c|}{0.22} & 0.15 & 0.42 \\ 
    \hhline{|--|~|-|-|-|-|-|-||-|}
    \multicolumn{2}{|c|}{HippoRAG 2} &  & \multicolumn{1}{c|}{\cellcolor{grannysmithapple} \textbf{0.80}} & \multicolumn{1}{c|}{0.77} & \multicolumn{1}{c|}{\cellcolor{grannysmithapple}\textbf{0.73}} & \multicolumn{1}{c|}{0.56} & \multicolumn{1}{c|}{0.29} & 0.28 & \cellcolor{grannysmithapple} \textbf{0.57} \\ 
    \hhline{==|~|======::=}
    \multicolumn{1}{|c|}{\multirow{4}{*}{\begin{tabular}[c]{@{}c@{}}PAI-1\\ (Ours)\end{tabular}}} & only Naive Retriever &  & \multicolumn{1}{c|}{0.68 / 0.63 / 0.56} & \multicolumn{1}{c|}{0.68 / 0.74 / 0.70} & \multicolumn{1}{c|}{0.65 / 0.62 / 0.50} & \multicolumn{1}{c|}{- / - / 0.29} & \multicolumn{1}{c|}{0.32 / 0.46 / 0.12} & - / - / 0.14 & 0.58 / 0.61 / 0.38 \\ 
    \hhline{|~|-|~|-|-|-|-|-|-||-|}
    \multicolumn{1}{|c|}{} & only Traversal Algorithms &  & \multicolumn{1}{c|}{0.68 / 0.61 / 0.55} & \multicolumn{1}{c|}{- / - / 0.67} & \multicolumn{1}{c|}{0.65 / 0.61 / 0.62} & \multicolumn{1}{c|}{0.45 / 0.35 / 0.34} & \multicolumn{1}{c|}{0.32 / 0.65 / 0.12} & 0.63 / 0.36 / \cellcolor{grannysmithapple} \textbf{0.35} & 0.54 / 0.51 / 0.44 \\ 
    \hhline{|~|-|~|-|-|-|-|-|-||-|}
    \multicolumn{1}{|c|}{} & Combined Retrieval (best) &  & \multicolumn{1}{c|}{\begin{tabular}[c]{@{}c@{}} - / - / 0.56 \\ BS+NR / E\end{tabular}} & \multicolumn{1}{c|}{\begin{tabular}[c]{@{}c@{}}0.66 / 0.82 / 0.73 \\ BS+NR / all\end{tabular}} & \multicolumn{1}{c|}{\begin{tabular}[c]{@{}c@{}}0.65 / 0.61 / 0.62 \\ BS+WC / all\end{tabular}} & \multicolumn{1}{c|}{\begin{tabular}[c]{@{}c@{}}0.46 / 0.50 / 0.40 \\ BS+NR / all\end{tabular}} & \multicolumn{1}{c|}{\begin{tabular}[c]{@{}c@{}}0.35 / 0.36 / 0.17 \\ BS+NR / all\end{tabular}} & \begin{tabular}[c]{@{}c@{}}0.63 / 0.36 / 0.35 \\ BS+WC / all\end{tabular} & 0.55 / 0.53 / 0.47 \\ 
    \hhline{|--|~|-|-|-|-|-|-||-|}
    \multicolumn{1}{|c|}{\multirow{7}{*}{\begin{tabular}[c]{@{}c@{}}PAI-2\\ (Ours)\end{tabular}}} & \begin{tabular}[c]{@{}c@{}}only Traversal Algorithms\\ + w/o plan enhancement\end{tabular} &  & \multicolumn{1}{c|}{- / - / 0.57} & \multicolumn{1}{c|}{0.74 / 0.88 / 0.71} & \multicolumn{1}{c|}{0.71 / 0.81 / 0.47} & \multicolumn{1}{c|}{0.44 / \textbf{0.84} / 0.28} & \multicolumn{1}{c|}{0.34 / 0.80 / 0.08} & \textbf{0.83} / \textbf{0.63} / 0.26 & 0.61 / 0.79 / 0.39 \\ 
    \hhline{|~|-|~|-|-|-|-|-|-||-|}
    \multicolumn{1}{|c|}{} & \begin{tabular}[c]{@{}c@{}}only Naive Retriever\\ + plan enhancement\end{tabular} &  & \multicolumn{1}{c|}{0.88 / \textbf{0.96} / 0.64} & \multicolumn{1}{c|}{0.89 / 0.88 / 0.77} & \multicolumn{1}{c|}{0.82 / 0.85 / 0.63} & \multicolumn{1}{c|}{0.60 / 0.78 / 0.48} & \multicolumn{1}{c|}{\textbf{0.68} / 0.86 / 0.28} & 0.66 / 0.60 / 0.26 & 0.75 / \textbf{0.82} / 0.51 \\ 
    \hhline{|~|-|~|-|-|-|-|-|-||-|}
    \multicolumn{1}{|c|}{} & \begin{tabular}[c]{@{}c@{}}only Traversal Algorithms\\ + plan enhancement\end{tabular} &  & \multicolumn{1}{c|}{0.91 / 0.93 / 0.67} & \multicolumn{1}{c|}{\textbf{0.92} / \textbf{0.91} / 0.77} & \multicolumn{1}{c|}{\textbf{0.89} / \textbf{0.87} / 0.67} & \multicolumn{1}{c|}{\textbf{0.74} / 0.80 / 0.54} & \multicolumn{1}{c|}{0.66 / \textbf{0.87} / 0.33} & \textbf{0.83} / 0.58 / 0.34 & \textbf{0.82} / \textbf{0.82} / 0.55 \\ 
    \hhline{|~|-|~|-|-|-|-|-|-||-|}
    \multicolumn{1}{|c|}{} & \begin{tabular}[c]{@{}c@{}}Combined Retrieval (best)\\ + plan enhancement\end{tabular} &  & \multicolumn{1}{c|}{\begin{tabular}[c]{@{}c@{}}\textbf{0.93} / 0.91 / 0.69\\ BS+NR / E\end{tabular}} & \multicolumn{1}{c|}{\cellcolor{grannysmithapple}\begin{tabular}[c]{@{}c@{}}0.89 / 0.83 / \textbf{0.8}\\ BS + NR / all\end{tabular}} & \multicolumn{1}{c|}{\begin{tabular}[c]{@{}c@{}}\textbf{0.89} / \textbf{0.87} / 0.67\\ BS+WC / all\end{tabular}} & \multicolumn{1}{c|}{\cellcolor{grannysmithapple}\begin{tabular}[c]{@{}c@{}}\textbf{0.74} / 0.74 / \textbf{0.58}\\ BS+NR / all\end{tabular}} & \multicolumn{1}{c|}{\cellcolor{grannysmithapple}\begin{tabular}[c]{@{}c@{}}0.66 / \textbf{0.87} / \textbf{0.33}\\ BS+WC / all\end{tabular}} & \begin{tabular}[c]{@{}c@{}}\textbf{0.83} / 0.57 / 0.34\\ BS+WC / E\end{tabular} & \cellcolor{grannysmithapple} \textbf{0.82} / 0.79 / \textbf{0.57} \\ 
    \hhline{=========::=}
    \multicolumn{3}{|c|}{Mean (by best setups)} & \multicolumn{1}{c|}{0.59} & \multicolumn{1}{c|}{0.68} & \multicolumn{1}{c|}{0.52} & \multicolumn{1}{c|}{0.38} & \multicolumn{1}{c|}{0.20} & 0.23 & 0.43 \\ 
    \hhline{|-|-|-|-|-|-|-|-|-||-|}
    \end{tabular}}
    \renewcommand{\arraystretch}{1}
    \caption{Best LLM-as-a-Judge scores for LightRAG, RAPTOR, HippoRAG 2 and Context Relevance / Faithfulness / LLM-as-a-Judge scores for PAI-1 and PAI-2 on six benchmarks. To prepare document stores and perform QA for every method Qwen2.5 7B was used. For PAI-1 and PAI-2 corresponding cells also contains retrieval algorithm and the type of restriction, applied to graph during traversal. Shortcuts for retrieval algorithms: "BS+WC" – hybrid of BeamSearch and WaterCircles; "BS+NR" – hybrid of BeamSearch and NaiverRetriever. Shortcuts for graph restrictions: "all" – no restrictions applied; "E" – episodic vertices were excluded from traversal.}
	\label{tab:main_comparison}
\end{table*}

%Из Таблицы~\ref{tab:main_comparison} видно, что PAI-2 превзошёл аналоги в среднем на $4\%$ по LLM-as-a-Judge на трёх из шести бенчмарках: TriviaQA, 2WikiMultihopQA и MuSiQue. При этом на HotpoQA и DiaASQ бенчмарках получен сопоставимый результат с HippoRAG и PAI-1 методами: с разницей $6\%$ и $1\%$ по LLM-as-a-Judge соответственно. Также можно увидеть, что самым лёгким бенчмарком (по сложности вопросов) оказался TriviaQA, а место самого сложного бенчмарка занял MuSiQue со средними показателями LLM-as-a-Judge $0.68$ и $0.20$ соответственно. В свою очередь, можно заметить серьёзное отставание PAI-2 от HippoRAG 2 на NaturalQuestions бенчмарке: с разницей $11\%$ по LLM-as-a-Judge. Это связано со следующими особенностями бенчмарка: (1) исходные вопросы представлены в нижнем регистре, что повышает вероятность пропуска важной именованной сущности и приводит к потере важной информации для генерации релевантного ответа и (2) некоторые вопросы требуют общих или недостаточно конкретизированных ответов, например, "how are the american declaration of independence and french declaration of the rights of man similar"; в виду того, что в PAI-2 нет явного механизма по детекции типа вопроса и ожидаемого формата ответа на этапе принятия решения LLM слабо уверена в полноте найденной/собранной информации и возвращает NoAnswer-заглушку. 

As shown in Table~\ref{tab:main_comparison}, PAI-2 achieves superior results with $4\%$ average gain by LLM-as-a-Judge on 3 out of 6 benchmarks: TriviaQA, 2WikiMultihopQA, and MuSiQue. Meanwhile, on HotpoQA and DiaASQ, it achieves comparable results to HippoRAG and PAI-1: with a $6\%$ and $1\%$ difference by LLM-as-a-Judge, respectively. It can also be seen that TriviaQA turned out to be the easiest benchmark (in terms of question difficulty), while MuSiQue took the place of the most difficult one, with average LLM-as-a-Judge scores $0.68$ and $0.20$, respectively. In turn, it can be noticed a significant gap between PAI-2 and HippoRAG 2 on the NaturalQuestions benchmark: with a $11\%$ difference by LLM-as-a-Judge. This observation may be attributed to the following characteristics of this benchmark: (1) the original questions are presented in lowercase, which increases probability to miss critical named entities and consequently lose essential information, required for relevant response generation; (2) some questions require general or insufficiently specific answers, for example, "how are the American declaration of independence and the French declaration of the rights of man similar"; because of PAI-2, that does not have an explicit mechanism for detecting question type and expected answer format, at the decision-making stage, it often returns "No Answer" response due to uncertainty, regarding the completeness of retrieved/summarized knowledge from memory.

%Мы провели серию экспериментов, в рамках которых оценили эффективность работы PAI-2 с выключенным механизмом корректировки плана поиска. Из Таблицы~\ref{tab:main_comparison} видно, что в сравнении с лучшими результатами PAI-2 качество генерируемых ответов упало на 18\% по LLM-as-a-Judge метрике. Это говорит о том, что некоторые вопросы требуют изменения стратегии поиска в зависимости о того набора информации, который хранится и доступен в моделе памяти. В качестве примера рассмотрим обработку вопроса "Do both films Payment On Demand and My Cousin From Warsaw have the directors from the same country?". Для него был сгенерирован следующий начальный план: (1) "Who is the director of the film Payment On Demand?", (2) "Who is the director of the film My Cousin From Warsaw?", (3) "What country is the director of Payment On Demand from?", (4) "What country is the director of My Cousin From Warsaw from?". Из него видно, что последние два шага требуют уточнения за счёт информации, полученной с первых, более конкретных и простых, шагов. Если этот план не корректировать, то для генерации финального ответа будет доступна следующая информация: (1) "Curtis Bernhardt", (2) "Carl Boese is the director of the film My Cousin From Warsaw.", (3) "Germany", (4) "<|NotEnoughtInfo|>" . Можно увидеть, что ее недостаточно для генерации релевантного ответа. В случае модификации шагов поиска на основе информации, полученной на промежуточных шагах по вопросу будет получен следующий план: (1) "Who is the director of the film Payment On Demand?", (2) "Who is the director of the film My Cousin From Warsaw?", (3) "What country is the director of Payment On Demand from?", (4) "What country is Carl Boese from?", (5) "What country is Curtis Bernhardt from?". На его основе для генерации финального ответа будет доступна следующая информация: (1) "The director of the film Payment On Demand is Curtis Bernhardt.", (2) "The director of the film "My Cousin from Warsaw" is Carl Boese.", (3) "Curtis Bernhardt, the director of Payment On Demand, was born in New York, New York.", (4) " Carl Boese was a German film director, screenwriter, and producer.", (5) "Curtis Bernhardt, born as Kurt Bernhardt in Worms, Germany, was from Germany." . Можно увидеть, что за счёт улучшения последних шагов удалось извлечь важную информацию для генерации точного ответа.

We conducted a series of experiments to evaluate PAI-2 with disabled search plan enhancement mechanism. Table~\ref{tab:main_comparison} shows that, compared to best PAI-2`s configurations, accuracy of generated answers is degraded on $18\%$ by LLM-as-a-Judge. This observation can be attributed to the form/complexity of some questions, that requires a dynamic search strategy with intermediate grounding on available knowledge in a memory graph. As an example, consider the question "Do both films Payment On Demand and My Cousin From Warsaw have the directors from the same country?". For it  the following initial plan was generated: (1) "Who is the director of the film Payment On Demand?"; (2) "Who is the director of the film My Cousin From Warsaw?"; (3) "What country is the director of Payment On Demand from?"; (4) "What country is the director of My Cousin From Warsaw from?". It can be seen that the last two steps require clarification with use of information, obtained from the first, more specific and clear, steps. With disabled plan enhancement, the following information will be available by each search step to generate final answer: (1) "Curtis Bernhardt"; (2) "Carl Boese is the director of the film My Cousin From Warsaw."; (3) "Germany"; (4) "<|NotEnoughtInfo|>" . It can be seen that it is not enough to generate a relevant answer. If the next search steps can be modified based on information, obtained by the previous ones, we get the following plan: (1) "Who is the director of the film Payment On Demand?"; (2) "Who is the director of the film My Cousin From Warsaw?"; (3) "What country is the director of Payment On Demand from?", (4) "What country is Carl Boese from?"; (5) "What country is Curtis Bernhardt from?". With that enhanced plan the following information will be available to generate the final answer: (1) "The director of the film Payment On Demand is Curtis Bernhardt."; (2) "The director of the film "My Cousin from Warsaw" is Carl Boese."; (3) "Curtis Bernhardt, the director of Payment On Demand, was born in New York, New York."; (4) "Carl Boese was a German film director, screenwriter, and producer."; (5) "Curtis Bernhardt, born as Kurt Bernhardt in Worms, Germany, was from Germany.". It can be seen that by improving the last steps, it was possible to extract crucial information from memory to generate an accurate answer.

%В сравнении с нашей предыдущей версией PAI фреймворка (PAI-1) предложенный алгоритм поиска (QA-пайплайн) в рамках PAI-2 позволил значительность повысить качество генерируемых ответов. Средний прирост значений Context Relevance, Faithfullness и LLM-as-a-Judge метрик составил 27\%, 26\% и 10\% соответственно. Т.е. добавление стадии планирования с механизмом корректировки поиска, а также выполнения обхода графа знаний на основе набора конкретизированных вопросов-подсказок повышает вероятность извлечения релевантной информации из структурированного массива документов, а также повышает  консистентность финального ответа с имеющимися набором знаний. В свою очередь, видна тенденция в превосходстве алгоритмов обхода графа над стандартным плоским ретривером: как для PAI-1, так и для PAI-2 наблюдается средний прирост в 5\% по LLM-as-a-Judge. При этом, важно контролировать количество шума в извлеченных триплетах, которое подаётся LLM для суммаризации накопленной информации и принятия решения о последующем шаге поиска. На примере использования NaiveRetriever алгоритма из Таблицы~\ref{tab:nr_vs_gt_comparison} видно, что исключение episodic триплетов из LLM контекста как для PAI-1, так и для PAI-2 позволяет смягчить влияние Lost in the Middle проблемы и повысить точность и обоснованность генерируемых ответов.

Compared to our previous version of the PAI framework (PAI-1), proposed QA pipeline in PAI-2 significantly improves answers quality. An average increase by Context Relevance, Faithfulness, and LLM-as-a-Judge metrics are 27\%, 26\% and 10\%, respectively. This means that integration of planning stage with a search steps enhancement mechanism and knowledge graph traversal based on a set of detailed/adjusted clue-queries increases probability to extract relevant information from a structured document store and improves consistency of final answer with existing knowledge base. Furthermore, a trend toward the superiority of graph traversal algorithms over the standard flattened retriever can be observed: both PAI-1 and PAI-2 demonstrate an average 5\% increase by LLM-as-a-Judge. Notably, it is important to control the amount of noise in the extracted triples, which LLM is using for knowledge summarization and decision making about the next search step. On example of using NaiveRetriever algorithm, in Table~\ref{tab:nr_vs_gt_comparison} it can be seen that exclusion of episodic triples from LLM context for both PAI-1 and PAI-2 mitigates the impact of "Lost in the Middle" problem~\cite{liu-etal-2024-lost} and improves accuracy and groundedness of answers.

\begin{table*}[th!]
    \centering
	\renewcommand{\arraystretch}{1.3}
	\resizebox{\textwidth}{!}{%
    \begin{tabular}{|c|c|c|c|c|c|c|c|c||c|}
    \hhline{|-|-|-|------||-}
    \multicolumn{1}{|c|}{\multirow{2}{*}{\textbf{Method}}} & \multicolumn{1}{c|}{\multirow{2}{*}{\begin{tabular}[c]{@{}c@{}}\textbf{Accepted} \\ \textbf{Triplet Types}\end{tabular}}} & \multirow{2}{*}{\textbf{LLM}} & \multicolumn{6}{c||}{\textbf{Dataset}} & \multirow{2}{*}{Mean} \\ 
    \hhline{|~|~|~|------||~|}
    \multicolumn{1}{|c|}{} & \multicolumn{1}{c|}{} &  & \textbf{Natural Questions} & \textbf{TriviaQA} & \textbf{HotpotQA} & \textbf{2WikiMultihopQA} & \textbf{MuSiQue} & \textbf{DiaASQ} &  \\ 
    \hhline{=========::=}
    \multicolumn{1}{|c|}{\multirow{4}{*}{PAI-1}} & \multicolumn{1}{c|}{simple} & \multirow{8}{*}{Qwen2.5 7B} & 0.63 / 0.62 / 0.47 & 0.63 / 0.66 / 0.68 & 0.63 / 0.45 / 0.48 & - / - / 0.26 & \cellcolor{grannysmithapple} 0.23 / 0.52 / \textbf{0.13} & \cellcolor{lightsalmonpink} - / - / 0.07 & 0.53 / 0.56 / 0.35 \\ 
    \hhline{|~|-|~|-|-|-|-|-|-||-|}
    \multicolumn{1}{|c|}{} & \multicolumn{1}{c|}{hyper} &  & 0.60 / 0.57 / 0.45 & 0.63 / 0.69 / 0.64 & - / - / 0.35 & - / - / 0.23 & \cellcolor{lightsalmonpink} 0.21 / \textbf{0.68} / 0.06 & \cellcolor{grannysmithapple} - / - / \textbf{0.15} & 0.48 / \textbf{0.65} / 0.31 \\ 
    \hhline{|~|-|~|-|-|-|-|-|-||-|}
    \multicolumn{1}{|c|}{} & \multicolumn{1}{c|}{simple, hyper} &  & \cellcolor{grannysmithapple} \textbf{0.68} / \textbf{0.64} / \textbf{0.56} & \cellcolor{grannysmithapple}\textbf{0.68} / \textbf{0.74} / \textbf{0.70} & \cellcolor{grannysmithapple} \textbf{0.65} / \textbf{0.62} / \textbf{0.50} & \cellcolor{grannysmithapple} - / - / \textbf{0.29} & \textbf{0.32} / 0.46 / 0.12 & - / - / 0.14 & \cellcolor{grannysmithapple} \textbf{0.58} / 0.62 / \textbf{0.38} \\ 
    \hhline{|~|-|~|-|-|-|-|-|-||-|}
    \multicolumn{1}{|c|}{} & \multicolumn{1}{c|}{episodic} &  &  \cellcolor{lightsalmonpink} 0.48 / 0.61 / 0.37 &  \cellcolor{lightsalmonpink} 0.65 / \textbf{0.74} / 0.44 & \cellcolor{lightsalmonpink} - / - / 0.21 & \cellcolor{lightsalmonpink} 0.41 / 0.51 / 0.13 & 0.16 / 0.56 / - & - / - / - & \cellcolor{lightsalmonpink} 0.42 / 0.60 / 0.29 \\ 
    \hhline{==|~|======::=}
    \multicolumn{1}{|c|}{\multirow{4}{*}{PAI-2}} & \multicolumn{1}{c|}{simple} &  & \cellcolor{lightsalmonpink} 0.86 / \textbf{0.96} / 0.58 & 0.86 / 0.87 / 0.69 & \textbf{0.83} / \textbf{0.85} / 0.59 & 0.59 / 0.67 / \textbf{0.48} & \cellcolor{lightsalmonpink} \textbf{0.69} / \textbf{0.86} / 0.25 & 0.49 / 0.56 / 0.10 & 0.72 / 0.80 / 0.45 \\ 
    \hhline{|~|-|~|-|-|-|-|-|-||-|}
    \multicolumn{1}{|c|}{} & \multicolumn{1}{c|}{hyper} &  & \textbf{0.90} / 0.92 / 0.59 & 0.88 / 0.86 / 0.72 & \cellcolor{lightsalmonpink} 0.73 / 0.79 / 0.50 & 0.52 / 0.74 / 0.34 & 0.59 / 0.73 / 0.27 & 0.62 / 0.61 / 0.25 & 0.71 / 0.78 / 0.44 \\ 
    \hhline{|~|-|~|-|-|-|-|-|-||-|}
    \multicolumn{1}{|c|}{} & \multicolumn{1}{c|}{simple, hyper} &  & 0.88 / \cellcolor{grannysmithapple}  \textbf{0.96} / \textbf{0.64} & \cellcolor{grannysmithapple} \textbf{0.89} / \textbf{0.88} / \textbf{0.77} & \cellcolor{grannysmithapple} 0.82 / \textbf{0.85} / \textbf{0.63} & \cellcolor{grannysmithapple} \textbf{0.60} / \textbf{0.78} / \textbf{0.48} & 0.68 / \textbf{0.86} / 0.28 & \cellcolor{grannysmithapple} \textbf{0.66} / 0.60 / \textbf{0.26} & \cellcolor{grannysmithapple} \textbf{0.76} / \textbf{0.82} / \textbf{0.51} \\ 
    \hhline{|~|-|~|-|-|-|-|-|-||-|}
    \multicolumn{1}{|c|}{} & \multicolumn{1}{c|}{episodic} &  & 0.75 / 0.93 / 0.60 & \cellcolor{lightsalmonpink} - / - / 0.66 & 0.73 / 0.81 / 0.55 & \cellcolor{lightsalmonpink} 0.48 / 0.75 / 0.33 & \cellcolor{grannysmithapple} 0.62 / 0.82 / \textbf{0.32} & \cellcolor{lightsalmonpink} 0.32 / \textbf{0.62} / 0.06 & \cellcolor{lightsalmonpink} 0.58 / 0.79 / 0.42 \\ 
    \hhline{=========::=}
    \multicolumn{3}{|c|}{Mean (PAI-1)} & 0.60 / 0.61 / 0.46 & 0.65 / 0.71 / 0.62 & 0.64 / 0.54 / 0.38 & 0.41 / 0.51 / 0.23 & 0.23 / 0.56 / 0.10 & - & 0.51 / 0.60 / 0.34 \\ 
    \hhline{|---|-|-|-|-|-|-||-|}
    \multicolumn{3}{|c|}{Mean (PAI-2)} & 0.85 / 0.94 / 0.60 & 0.88 / 0.87 / 0.71 & 0.78 / 0.82 / 0.57 & 0.55 / 0.74 / 0.41 & 0.64 / 0.82 / 0.28 & 0.52 / 0.60 / 0.17 & 0.70 / 0.79 / 0.46 \\
    \hhline{|---|-|-|-|-|-|-||-|}
    \end{tabular}}
    \renewcommand{\arraystretch}{1}
    \caption{PAI-1 and PAI-2 performance depending on accepted triples types for final answer generation across six datasets. For graph/triples traversal/retrieval NaiveRetriever algorithm is used}
	\label{tab:nr_vs_gt_comparison}
\end{table*}

To measure the impact of clue-queries number per search step on relevance of generated answers, a corresponding series of experiments was conducted: see Table~\ref{tab:main_maxcluequeries_ablation}. Our non aggregated results for that table are presented in Appendix~\ref{app:nonaggr_clueq}.

%Для оценки влияния количества генерируемых clue-вопросов у PAI-2 в рамках одного шага поиска на релевантность генерируемых ответов была проведена соответствующая серия ablation экспериментов: см. Таблицу~\ref{tab:nr_vs_gt_comparison}. Из Таблицы видно, что при увеличении числа clue-вопросов, на основе которых выполняется обход графа знаний, повышается качество ответов: получен прирост на 4\% по LLM-as-a-Judge метрике при изменении максимального количество clue-вопросов с 1 до 8. Данное наблюдение связано с алгоритмом формирования пространства памяти. При добавлении новых документов в модель памяти набор извлечённых триплетов проверяется на дубликаты с триплетами в существующем графе знаний. Данная проверка выполняется по прямому сопоставлению и сравнению текстовых атрибутов в триплетах. В виду того, что одно и то же знание может быть по разному сформулировано, в строящемся графе могут появлятся подграфы, которые содержат и описывают одно и тоже знание, но с помощью разных сущностей. При этом такие подграфы могут не иметь общих вершин и по отдельности быть неполными: разрозненно по подграфам содержать информацию об одном обьекте. В связи с этим, использование нескольких вариантов вершин для одной сущности из плана поиска и последующая генерация clue-вопросов на основе их линейной комбинации позволяет обойти такие подграфы, найти и собрать воедино запрашиваему информацию по заданному обьекту для генерации точного и полногл ответа.

\begin{table*}[th!]
    \centering
	\renewcommand{\arraystretch}{1.3}
	\resizebox{\textwidth}{!}{
    \begin{tabular}{|cc|c|c|c|c|c|c||c|}
    \hhline{|-|-|------||-|}
    \multicolumn{1}{|c|}{\multirow{2}{*}{\textbf{Max Clue Queries}}} & \multirow{2}{*}{\textbf{LLM}} & \multicolumn{6}{c||}{\textbf{Dataset}} & \multirow{2}{*}{Mean} \\ 
    \hhline{|~|~|------||~|}
    \multicolumn{1}{|c|}{} &  & \multicolumn{1}{c|}{\textbf{Natural Questions}} & \multicolumn{1}{c|}{\textbf{TriviaQA}} & \multicolumn{1}{c|}{\textbf{HotpotQA}} & \multicolumn{1}{c|}{\textbf{2WikiMultihopQA}} & \multicolumn{1}{c|}{\textbf{MuSiQue}} & \textbf{DiaASQ} &  \\ 
    \hhline{========::=}
    \multicolumn{1}{|c|}{1} & \multirow{5}{*}{Qwen2.5 7B} & \multicolumn{1}{c|}{\cellcolor{lightsalmonpink} 0.85 / \textbf{0.94} / 0.60} & \multicolumn{1}{c|}{\cellcolor{lightsalmonpink} 0.88 / \textbf{0.89} / 0.72} & \multicolumn{1}{c|}{\cellcolor{lightsalmonpink} 0.83 / \textbf{0.84} / 0.58} & \multicolumn{1}{c|}{\cellcolor{lightsalmonpink} 0.66 / 0.75 / 0.49} & \multicolumn{1}{c|}{\cellcolor{lightsalmonpink} 0.63 / 0.82 / 0.24} & \cellcolor{lightsalmonpink} 0.76 / \textbf{0.66} / 0.25 & \cellcolor{lightsalmonpink} 0.77 / \textbf{0.82} / 0.48 \\ 
    \hhline{|-|~|-|-|-|-|-|-||-|}
    \multicolumn{1}{|c|}{2} &  & \multicolumn{1}{c|}{\cellcolor{grannysmithapple} 0.89 / 0.93 / \textbf{0.66}} & \multicolumn{1}{c|}{0.88 / 0.86 / 0.73} & \multicolumn{1}{c|}{0.84 / \textbf{0.84} / 0.58} & \multicolumn{1}{c|}{0.68 / 0.75 / 0.51} & \multicolumn{1}{c|}{\cellcolor{grannysmithapple} 0.67 / \textbf{0.83} / \textbf{0.29}} & 0.80 / 0.64 / 0.26 & 0.79 / 0.81 / 0.50 \\ 
    \hhline{|-|~|-|-|-|-|-|-||-|}
    \multicolumn{1}{|c|}{4} &  & \multicolumn{1}{c|}{\textbf{0.90} / 0.92 / 0.65} & \multicolumn{1}{c|}{0.89 / 0.87 / 0.75} & \multicolumn{1}{c|}{0.88 / \textbf{0.84} / 0.61} & \multicolumn{1}{c|}{0.71 / 0.76 / 0.53} & \multicolumn{1}{c|}{0.69 / 0.82 / 0.25} & 0.82 / 0.65 / 0.27 & \textbf{0.82} / 0.81 / 0.51 \\ 
    \hhline{|-|~|-|-|-|-|-|-||-|}
    \multicolumn{1}{|c|}{6} &  & \multicolumn{1}{c|}{\textbf{0.90} / 0.93 / 0.64} & \multicolumn{1}{c|}{0.89 / 0.87 / \textbf{0.77}} & \multicolumn{1}{c|}{0.88 / \textbf{0.84} / \textbf{0.62}} & \multicolumn{1}{c|}{\cellcolor{grannysmithapple} \textbf{0.73} / \textbf{0.78} / \textbf{0.55}} & \multicolumn{1}{c|}{0.70 / 0.81 / 0.25} & \textbf{0.84} / 0.64 / 0.28 & \textbf{0.82} / 0.81 / \textbf{0.52} \\ 
    \hhline{|-|~|-|-|-|-|-|-||-|}
    \multicolumn{1}{|c|}{8} &  & \multicolumn{1}{c|}{\textbf{0.90} / 0.92 / 0.65} & \multicolumn{1}{c|}{\cellcolor{grannysmithapple} \textbf{0.90} / 0.88 / \textbf{0.77}} & \multicolumn{1}{c|}{\cellcolor{grannysmithapple} \textbf{0.89} / 0.83 / \textbf{0.62}} & \multicolumn{1}{c|}{0.71 / 0.76 / 0.53} & \multicolumn{1}{c|}{\textbf{0.71} / 0.82 / 0.25} & \cellcolor{grannysmithapple} \textbf{0.84} / 0.64 / \textbf{0.29} & \cellcolor{grannysmithapple} \textbf{0.82} / 0.81 / \textbf{0.52} \\ 
    \hhline{========::=}
    \multicolumn{2}{|c|}{Mean} & \multicolumn{1}{c|}{0.89 / 0.93 / 0.64} & \multicolumn{1}{c|}{0.89 / 0.87 / 0.75} & \multicolumn{1}{c|}{0.86 / 0.84 / 0.60} & \multicolumn{1}{c|}{0.70 / 0.76 / 0.52} & \multicolumn{1}{c|}{0.68 / 0.82 / 0.26} & 0.81 / 0.65 / 0.27 & 0.80 / 0.81 / 0.51 \\ 
    \hhline{|--|-|-|-|-|-|-||-|}
    \end{tabular}}
    \renewcommand{\arraystretch}{1}
    \caption{PAI-2`s QA pipeline performance depending on generated clue queries amount for each step of a search plan. Cells contain Context Relevance, Faithfulness and LLM-as-a-Judge scores.}
	\label{tab:main_maxcluequeries_ablation}
\end{table*}

From Table~\ref{tab:main_maxcluequeries_ablation} it can be seen that with increase of clue-queries number (using for manage graph traversal from associated starting vertices) answers quality is improves: an average 4\% gain by LLM-as-a-Judge was achieved when changing the maximum number of clue-queries from 1 to 8. This observation may be attributed to characteristics of PAI`s memory graph construction algorithm. When a new document is adding to memory, its extracted triplets are validating for duplicates with triplets in the existing memory graph. This validation and subsequent filtering (for duplicates) is performed by exact match metric for triplet`s text attributes. Because the same knowledge can be formulated in different ways, several subgraphs may appear in memory that contain and describe the same knowledge but using different entities. Such subgraphs may not share any common vertices and may be incomplete: information about a single object can be scattered across subgraphs. Therefore, using multiple vertices for a single entity from the search plan and subsequently generating clue questions based on their linear combination allows us to traverse such subgraphs, find and aggregate the requested information on a given object to generate an accurate and complete answer.

It is also important to note the required time to process a single user question with our method: see Table~\ref{tab:weak_medium_pai_latency}.

\begin{table*}[th!]
    \centering
	\renewcommand{\arraystretch}{1.3}
	%\resizebox{\textwidth}{!}{%
    \begin{tabular}{|cc|c|c|c|c|c|c||c|}
    \hhline{|-|-|------||-|}
    \multicolumn{1}{|c|}{\multirow{2}{*}{\textbf{Method}}} & \multirow{2}{*}{\textbf{LLM}} & \multicolumn{6}{c||}{\textbf{Dataset}} & \multirow{2}{*}{Mean} \\ 
    \hhline{|~|~|------||~|}
    \multicolumn{1}{|c|}{} &  & \textbf{Natural Questions} & \textbf{TriviaQA} & \textbf{HotpotQA} & \textbf{2WikiMultihopQA} & \textbf{MuSiQue} & \textbf{DiaASQ} &  \\ 
    \hhline{========::=}
    \multicolumn{1}{|c|}{PAI-1} & \multirow{2}{*}{Qwen2.5 7B} & 0.43 & 0.88 & 1.11 & 1.51 & 0.50 & 1.55 & 1.0 \\
    \hhline{|-|~|-|-|-|-|-|-||-|}
    \multicolumn{1}{|c|}{PAI-2} &  & 0.72 & 1.44 & 1.42 & 1.06 & 0.73 & 3.70 & 1.51 \\ 
    \hhline{========::=}
    \multicolumn{2}{|c|}{Mean} & 0.57 & 1.16 & 1.27 & 1.28 & 0.62 & 2.62 & 1.25  \\
    \hhline{|--|-|-|-|-|-|-||-|}
    \end{tabular}%}
    \renewcommand{\arraystretch}{1}
    \caption{Latency (in minutes) of PAI-1 and PAI-2 QA pipelines with Qwen2.5 7B across 6 datasets}
	\label{tab:weak_medium_pai_latency}
\end{table*}

%Table~\ref{tab:weak_medium_pai_latency} shows that PAI-2 requires, approximately, double time to process one question and generated an answer compared to PAI-1. Это связано с тем, что в PAI-1 выполняется только одна итерация поиска информации в графе, когда в PAI-2 количество итераций может варьироваться в зависимости от сложности вопроса. Бутолочным горлышком по части производительности являются следующие операции: (1) инференс LLM, (2) векторный поиск и (3) обход графа знаний. Чтобы смягчить влияние данных факторов на скорость работы поискового алгоритма необходимо прибегать к следующим практикам: (1) кешировать и переиспользовать результаты инференса LLM, (2) разбивать большие векторные хранилища на небольшие подмножетсва за счёт типизации объектов и (3) ограничение пространства поиска в графе. 

Table~\ref{tab:weak_medium_pai_latency} shows that PAI-2 requires, approximately, double time to process one question and generate an answer compared to PAI-1. This is due to the fact that PAI-1 performs only a single iteration of information retrieval, while in PAI-2 the number of iterations can vary depending on the complexity of the question. The following operations can be stated as bottlenecks: (1) LLM inference; (2) vector search; (3) knowledge graph traversal. To mitigate the impact of these factors on performance of the search workflow, it is necessary to use the following practices: (1) caching and reusing LLM inference results; (2) split large vector stores into smaller subsets by elements types; (3) limit the search space in memory graph.

Additionally, we evaluate our plain-text-to-knowledge-graph extraction algorithm (Memorize pipeline) on MINE benchmark to measure the factual completeness of constructing PAI`s memory graphs. Our method demonstrates SOTA results with 89\% information-retention score: more details are provided in Appendix~\ref{app:mine_comparison}. 

%\subsection{Error Analysis}
% https://arxiv.org/pdf/2409.03155
% круговая диаграмма
% TODO